
\documentclass[11pt]{article}
\usepackage[margin=1in]{geometry}
\usepackage[T1]{fontenc}
\usepackage{lmodern}
\usepackage{microtype}
\usepackage{amsmath,amssymb,amsthm,mathtools}
\usepackage{booktabs,longtable,array}
\usepackage{enumitem}
\usepackage{xcolor}
\usepackage{hyperref}
\usepackage[nameinlink,capitalise,noabbrev]{cleveref}
\usepackage{listings}
\usepackage{graphicx}
\usepackage{fancyhdr}
\usepackage{setspace}
\usepackage{bm}
\usepackage{url}

\definecolor{deepblue}{RGB}{25,57,106}
\definecolor{deepred}{RGB}{135,35,35}
\definecolor{softgray}{RGB}{245,246,248}
\hypersetup{colorlinks=true,linkcolor=deepblue,citecolor=deepblue,urlcolor=deepblue,pdfauthor={OpenAI GPT-5.6 Pro},pdftitle={Local factorials, dyadic capacity, and synchronization barriers for the Alaoglu--Erdos problem}}
\pagestyle{fancy}
\fancyhf{}
\fancyhead[L]{Alaoglu--Erd\H{o}s progress report}
\fancyhead[R]{\thepage}
\setlength{\headheight}{14pt}
\setlist[itemize]{leftmargin=2em,itemsep=0.25em,topsep=0.35em}
\setlist[enumerate]{leftmargin=2.2em,itemsep=0.25em,topsep=0.35em}
\lstset{basicstyle=\ttfamily\small,backgroundcolor=\color{softgray},frame=single,breaklines=true,columns=fullflexible,keepspaces=true}

\newtheorem{theorem}{Theorem}[section]
\newtheorem{proposition}[theorem]{Proposition}
\newtheorem{lemma}[theorem]{Lemma}
\newtheorem{corollary}[theorem]{Corollary}
\newtheorem{conjecture}[theorem]{Conjecture}
\newtheorem{criterion}[theorem]{Criterion}
\theoremstyle{definition}
\newtheorem{definition}[theorem]{Definition}
\newtheorem{problem}[theorem]{Problem}
\newtheorem{experiment}[theorem]{Experiment}
\theoremstyle{remark}
\newtheorem{remark}[theorem]{Remark}
\newtheorem{warning}[theorem]{Warning}

\newcommand{\N}{\mathbb N}
\newcommand{\Z}{\mathbb Z}
\newcommand{\Q}{\mathbb Q}
\newcommand{\R}{\mathbb R}
\newcommand{\Zp}{\mathbb Z_p}
\newcommand{\vp}{v_p}
\newcommand{\ord}{\operatorname{ord}}
\newcommand{\lcm}{\operatorname{lcm}}
\newcommand{\Int}{\operatorname{Int}}
\newcommand{\Vdm}{\operatorname{Vdm}}
\newcommand{\covol}{\operatorname{covol}}
\newcommand{\Span}{\operatorname{span}}
\newcommand{\calS}{\mathcal S}
\newcommand{\calE}{\mathcal E}
\newcommand{\calC}{\mathcal C}
\newcommand{\calL}{\mathcal L}
\newcommand{\floor}[1]{\left\lfloor #1\right\rfloor}
\newcommand{\abs}[1]{\left|#1\right|}
\newcommand{\norm}[1]{\left\lVert#1\right\rVert}

\title{\textbf{Local Factorials, Dyadic Capacity, and the Synchronization Barrier}\\[0.4em]
\large A nonduplicative progress report on the Alaoglu--Erd\H{o}s exponent problem}
\author{OpenAI GPT-5.6 Pro}
\date{23 August 2026}

\begin{document}
\maketitle

\begin{abstract}
Let
\[
   \calS=\{2^a3^b:a,b\in\N_0\}.
\]
The Alaoglu--Erd\H{o}s problem asks whether a real number $x$ satisfying
$2^x,3^x\in\Z$ must be an integer.  The attached report already develops a broad collection of transcendence, approximation, determinant, $p$-adic, and computational ideas.  The present report deliberately does not repeat that material.  It develops one interface that was previously only prospective: Bhargava $p$-orderings and generalized factorials of the $3$-smooth semigroup.

The main unconditional results are these.  First, the complete local factorial profile of $\calS$ is reduced to explicit formulas.  At $p=3$ it is the ordinary factorial profile.  At every odd prime $p\ge5$ it is governed by the order of the subgroup generated by $2$ and $3$ modulo $p$, together with one lifting index.  At $p=2$ there is an exact self-similar residue-tree recurrence.  Its Robin constant, or asymptotic insertion cost, is
\[
        V_2(\calS)=\frac{\sqrt{11}-1}{2}.
\]
Second, combining the odd-prime formula with the theorem of Bugeaud--Corvaja--Zannier on
$\gcd(2^d-1,3^d-1)$ gives
\[
        \log(n!_{\calS})=o(n^2).
\]
This subquadratic estimate is exactly the size needed by a high-node Newton interpolation attack.

The same analysis also uncovers a decisive obstruction.  A $2$-adic optimal ordering asymptotically places
$(\sqrt{11}-1)/5$ of its nodes in the odd branch, whereas a $3$-adic optimal ordering places $2/3$ of its nodes among the $3$-adic units.  Away from the node $1$, those two classes are disjoint, but their target masses sum to more than one.  Consequently no family of high simultaneous $p$-ordering prefixes can exist.  A quantitative version gives an explicit positive quadratic synchronization defect.  Thus the simplest scalar Newton-coefficient strategy cannot prove the conjecture, even though every one of its local estimates is favorable.

The report finishes by replacing that failed scalar target with a concrete rank-two coefficient-lattice program.  The proposed object is computable by Smith normal form, has a rigorous arithmetic lower bound, and admits a generalized-Wronskian archimedean estimate.  No proof of the Alaoglu--Erd\H{o}s conjecture is claimed here; the advance is an exact local theory, a global subquadratic theorem, a rigorous no-go result for the most tempting interpolation route, and a sharply specified next lemma whose truth can be tested computationally and formalized in Lean.
\end{abstract}

\tableofcontents
\newpage

\section{Executive summary and status}

\subsection{The mathematical status of this report}
The target implication is
\begin{equation}\label{eq:AE}
        2^x,3^x\in\Z \quad\Longrightarrow\quad x\in\Z,
        \qquad x\in\R.
\end{equation}
This report does not claim a proof of \eqref{eq:AE}.  It does prove several new-to-this-project structural theorems and, equally importantly, rules out a natural proof architecture that otherwise looks compelling when examined prime by prime.

The genuinely new work is concentrated in five statements.
\begin{enumerate}
\item An exact formula for the Bhargava $p$-sequence of $\calS$ at every odd prime $p\ge5$; see \cref{thm:odd-profile}.
\item An exact dyadic residue-tree recurrence and the closed form
$V_2=(\sqrt{11}-1)/2$; see \cref{thm:dyadic-recurrence,thm:dyadic-capacity}.
\item The global estimate $\log(n!_{\calS})=o(n^2)$; see \cref{thm:subquadratic}.
\item A first exact synchronization failure: the natural simultaneous Newton chain
$1,2,3,4,6$ cannot be extended by a sixth node, and the nearest continuation incurs the new prime factor $7$; see \cref{prop:sixth-node}.
\item A two-prime mass incompatibility and a quantitative quadratic energy gap; see \cref{thm:no-high-simultaneous,thm:quadratic-gap}.
\end{enumerate}

\subsection{Why the result is useful despite the obstruction}
The subquadratic factorial theorem proves that the local arithmetic is not too large.  If one could find simultaneous high Newton nodes $a_0,\dots,a_n\in\calS$ with $\log\min a_i\gg n$, then the generalized mean-value theorem would make the $n$th divided difference of $t^x$ smaller than the reciprocal of its forced denominator, and \eqref{eq:AE} would follow.  This conditional argument is completely rigorous and is recorded in \cref{thm:scalar-bridge}.

The obstruction theorem then explains why attempts to construct such a sequence keep failing: the locally optimal $2$-adic and $3$-adic node populations cannot coexist at high archimedean height.  This is not a vague failure of a search heuristic.  It is an exact incompatibility of equilibrium masses, with a positive numerical gap
\[
   \eta=\frac{\sqrt{11}-1}{5}-\frac13
       =\frac{3\sqrt{11}-8}{15}>0.
\]
The right response is therefore not a larger search for one universal chain.  It is to raise the arithmetic rank of the interpolation object.  The rank-two coefficient lattice in \cref{sec:rank-two} is designed precisely to absorb the two incompatible exceptional branches.

\subsection{Independence from competitive claims}
Recent claims about automated mathematical discoveries are not used as premises anywhere below.  Every theorem in this report is either proved from standard published results, reduced to an explicit finite recurrence, or clearly labeled as a conjectural next step.  This separation is essential: a competitive deadline is a reason to search aggressively, not a reason to lower the standard of proof.

\section{Normalization of the exponent problem}

\subsection{The semigroup formulation}
Assume that $x\in\R$ and put
\[
       m=2^x,\qquad n=3^x.
\]
If $m,n\in\Z$, then $m,n>0$ and, for every $a,b\ge0$,
\begin{equation}\label{eq:all-smooth-values}
       (2^a3^b)^x=m^a n^b\in\Z.
\end{equation}
Thus the problem is equivalent to classifying the real exponents for which the power function
\[
       f_x(t)=t^x
\]
is integer-valued on the multiplicative semigroup $\calS$.

This elementary reformulation is the reason integer-valued polynomial methods enter naturally.  We do not know $f_x$ on all positive integers, but we know it on an infinite, rank-two, multiplicatively generated set.  Its counting function is
\[
  \#\{s\in\calS:s\le X\}
   =\frac{(\log X)^2}{2\log2\log3}+O(\log X).
\]
The set is sparse additively, but its exponent lattice is two-dimensional.

\subsection{Elementary boundary reductions}
\begin{proposition}\label{prop:boundary}
Suppose $2^x,3^x\in\Z$.
Then $x\ge0$.  If $x\in\Q$, then $x\in\Z$.  If $x$ is algebraic, then $x\in\Z$.  Consequently every hypothetical counterexample is a positive transcendental number.
\end{proposition}
\begin{proof}
If $x<0$, then $0<2^x<1$, so $2^x$ is not an integer.  Write a rational $x=r/q$ in lowest terms, with $q>0$.  If $2^{r/q}$ is an integer, then its $q$th power is $2^r$; unique factorization forces $q\mid r$, hence $q=1$.  Finally, if $x$ is algebraic irrational, the Gelfond--Schneider theorem makes $2^x$ transcendental, contradicting its integrality.
\end{proof}

\begin{remark}
The four exponentials conjecture would settle the problem immediately: the two rows $\log2,\log3$ are linearly independent over $\Q$, and a nonrational $x$ makes the two columns $1,x$ linearly independent.  All four exponentials $2,3,2^x,3^x$ would nevertheless be algebraic.  The purpose of the present report is to exploit the additional fact that the last two exponentials are integers and that all values in \eqref{eq:all-smooth-values} are integers.
\end{remark}

\section{Bhargava orderings and generalized factorials}

\subsection{Local insertion costs}
Fix a prime $p$.  For distinct $a_0,\dots,a_{r-1}\in\calS$ and a candidate $a\in\calS$, define the insertion cost
\[
       C_p(a;a_0,\dots,a_{r-1})
       =v_p\!\left(\prod_{i<r}(a-a_i)\right)
       =\sum_{k\ge1}\#\{i<r:a_i\equiv a\pmod{p^k}\}.
\]
A $p$-ordering is a sequence $a_0,a_1,\dots$ in which $a_r$ minimizes this cost at every stage.  Bhargava proved that the minimum exponent
\[
       w_p(r)=C_p(a_r;a_0,\dots,a_{r-1})
\]
is independent of the choices made within ties.  The generalized factorial of $\calS$ is
\begin{equation}\label{eq:generalized-factorial}
       r!_{\calS}=\prod_p p^{w_p(r)}.
\end{equation}
For a fixed $r$ only finitely many factors are nontrivial.

The residue-tree identity in the definition is useful enough to emphasize.  The valuation of a product of differences is the sum, over every depth $k$, of the number of previously selected points in the same occupied residue class modulo $p^k$.  Thus a $p$-ordering is a greedy load-balancing process on the rooted tree of occupied $p$-adic residue classes.

\subsection{Dependence only on the $p$-adic closure}
Let $E_p$ be the closure of $\calS$ in $\Z_p$.  Every finite insertion cost is determined modulo a sufficiently high power of $p$, so the sequence $w_p(r)$ depends only on $E_p$.  This allows us to replace the discrete semigroup by a compact residue tree.  At odd primes not dividing $6$, the closure is a compact multiplicative subgroup.  At $p=3$ it is all of $\Z_3$.  At $p=2$ it is a self-similar union of one odd group and a scaled copy of itself.

\subsection{Pair energy}
For later use, define the minimum pair energy of $N$ points in a compact $p$-adic set $E$ by
\[
  F_E(N)=\min_{x_1,\dots,x_N\in E\atop x_i\ne x_j}
       \sum_{1\le i<j\le N}v_p(x_i-x_j),
       \qquad F_E(0)=F_E(1)=0.
\]
The residue-tree decomposition gives
\[
  \sum_{i<j}v_p(x_i-x_j)
  =\sum_{k\ge1}\sum_{c\bmod p^k}\binom{N_{k,c}}2,
\]
where $N_{k,c}$ is the occupancy of residue class $c$.  Convexity of $u\mapsto\binom u2$ and induction down the residue tree show that a greedy $p$-ordering gives nested energy minimizers.  Hence
\begin{equation}\label{eq:energy-increments}
       F_{E_p}(N)=\sum_{r=0}^{N-1}w_p(r),
       \qquad w_p(r)=F_{E_p}(r+1)-F_{E_p}(r).
\end{equation}
This identity makes self-similar recurrences particularly transparent.

\section{The complete odd-prime profile}

\subsection{The prime $3$}
\begin{theorem}\label{thm:p3}
The semigroup $\calS$ is dense in $\Z_3$.  Consequently
\[
        w_3(r)=v_3(r!)=\sum_{j\ge1}\floor{\frac r{3^j}}.
\]
\end{theorem}
\begin{proof}
The integer $2$ is a primitive root modulo $3^k$ for every $k\ge1$.  Let $y$ be a residue modulo $3^k$.  If $y\equiv0$, choose any $3^b$ with $b\ge k$.  Otherwise write $y=3^b u$ with $0\le b<k$ and $u$ a unit modulo $3^{k-b}$.  Choose $a$ with $2^a\equiv u\pmod{3^{k-b}}$.  Then $2^a3^b\equiv y\pmod{3^k}$.  Thus every residue class at every depth is occupied.  The ordinary ordering $0,1,2,\dots$ of $\Z_3$ has insertion cost $v_3(r!)$, and density transfers that profile to $\calS$.
\end{proof}

\subsection{Odd primes not dividing $6$}
For $p\ge5$, set
\[
 H_p=\overline{\langle2,3\rangle}\subset\Z_p^\times,
 \qquad r_p(k)=\#(H_p\bmod p^k).
\]
The nonnegative semigroup generated by $2$ and $3$ is dense in the closed group generated by them: in a compact group, inverses are limits of positive powers.  Since $\Z_p^\times$ is procyclic for odd $p$, so is $H_p$.

\begin{lemma}\label{lem:procyclic-profile}
Let $H\subset\Z_p^\times$ be a procyclic compact group and let
$r(k)=\#(H\bmod p^k)$.  Then its $p$-sequence is
\[
       w_H(n)=\sum_{k\ge1}\floor{\frac n{r(k)}}.
\]
\end{lemma}
\begin{proof}
Choose a topological generator $g$.  Among $1,g,\dots,g^{n-1}$, every residue class modulo $p^k$ occurs either $\floor{n/r(k)}$ or $\lceil{n/r(k)}\rceil$ times.  The candidate $g^n$ lies in a class occurring exactly $\floor{n/r(k)}$ times.  Therefore its total insertion cost is the displayed sum.  Every other candidate lies, at each depth, in a class with at least that occupancy.  Hence the sequence is a $p$-ordering.
\end{proof}

Define
\begin{equation}\label{eq:dp-cp}
 d_p=\lcm\bigl(\ord_p(2),\ord_p(3)\bigr),
 \qquad
 c_p=\min\bigl(v_p(2^{d_p}-1),v_p(3^{d_p}-1)\bigr).
\end{equation}
Here $d_p\mid p-1$ and $c_p\ge1$.

\begin{theorem}[Exact odd-prime profile]\label{thm:odd-profile}
For every prime $p\ge5$ and $k\ge1$,
\[
       r_p(k)=d_p p^{\max(0,k-c_p)}.
\]
Consequently
\begin{equation}\label{eq:odd-wp}
 w_p(n)=c_p\floor{\frac n{d_p}}
       +\sum_{j\ge1}\floor{\frac n{d_p p^j}}.
\end{equation}
In particular,
\[
       \lim_{n\to\infty}\frac{w_p(n)}n
       =\frac1{d_p}\left(c_p+\frac1{p-1}\right).
\]
\end{theorem}
\begin{proof}
The reduction of $H_p$ modulo $p$ has order $d_p$.  Its kernel in $1+p\Z_p$ is generated topologically by $2^{d_p}$ and $3^{d_p}$.  The $p$-adic logarithm identifies $1+p\Z_p$ with $p\Z_p$.  For $u\in1+p\Z_p$, one has $v_p(\log u)=v_p(u-1)$.  Hence the additive closed subgroup generated by
$\log(2^{d_p})$ and $\log(3^{d_p})$ is $p^{c_p}\Z_p$.  Its image modulo $p^k$ has size $p^{\max(0,k-c_p)}$.  Multiplying by the torsion factor $d_p$ proves the formula for $r_p(k)$.  Substitution into \cref{lem:procyclic-profile} gives \eqref{eq:odd-wp}; summing the geometric tail gives the limit.
\end{proof}

\begin{remark}
Usually $c_p=1$.  A value $c_p>1$ requires both $2^{d_p}-1$ and $3^{d_p}-1$ to have unexpectedly high $p$-adic valuation.  The formula isolates this joint lifting anomaly exactly; no heuristic assumption about Wieferich primes is needed.
\end{remark}

\section{The dyadic residue tree}

\subsection{The odd branch}
Let
\[
       H=\overline{\{3^b:b\ge0\}}\subset\Z_2^\times.
\]
The order of $3$ modulo $2^k$ is $1$ for $k=1$, $2$ for $k=2$, and $2^{k-2}$ for $k\ge3$.  Equivalently,
\[
       H=(1+8\Z_2)\;\sqcup\;(3+8\Z_2).
\]
Applying \cref{lem:procyclic-profile} gives the exact insertion cost in the odd branch:
\begin{equation}\label{eq:h2}
       h(n):=w_H(n)
       =n+\floor{\frac n2}+v_2(n!).
\end{equation}
Its asymptotic slope is $5/2$.

\subsection{Self-similarity of the full closure}
Let
\[
       E=\overline{\calS}^{\,2}\subset\Z_2.
\]
Every element is either odd and in $H$, or even and twice another element of $E$.  Thus
\begin{equation}\label{eq:E-decomp}
       E=H\sqcup 2E.
\end{equation}
The two components are in distinct residue classes modulo $2$, so cross-differences have valuation zero.

\begin{theorem}[Exact dyadic recurrence]\label{thm:dyadic-recurrence}
Let $F(N)=F_E(N)$ and
\[
       F_H(N)=\sum_{j=0}^{N-1}\left(j+\floor{\frac j2}+v_2(j!)\right).
\]
Then $F(0)=F(1)=0$ and, for $N\ge2$,
\begin{equation}\label{eq:F-recurrence}
       F(N)=\min_{0\le K\le N-1}
       \left\{F(K)+\binom K2+F_H(N-K)\right\}.
\end{equation}
The dyadic factorial exponent is
\[
       w_2(n)=F(n+1)-F(n).
\]
\end{theorem}
\begin{proof}
Take an $N$-point configuration and suppose $K$ points lie in $2E$.  Dividing those points by $2$ leaves a configuration in $E$, while each of its $\binom K2$ pair differences loses exactly one unit of valuation.  The remaining $N-K$ points lie in $H$, and all cross-valuations are zero.  Hence the minimum with a fixed $K$ is
$F(K)+\binom K2+F_H(N-K)$.  An energy minimizer with $N\ge2$ cannot have $K=N$, because dividing every point by $2$ would lower the energy by $\binom N2$.  This yields \eqref{eq:F-recurrence}.  The last assertion follows from \eqref{eq:energy-increments}.
\end{proof}

The recurrence is finite, integer-valued, and has no search over semigroup elements.  The first values are displayed below; each pair of rows gives $n$ and $w_2(n)$.
\begin{center}
\small
\begin{tabular}{*{10}{r}}
0 & 1 & 2 & 3 & 4 & 5 & 6 & 7 & 8 & 9 \\
0 & 0 & 1 & 1 & 3 & 4 & 4 & 5 & 7 & 9 \\
10 & 11 & 12 & 13 & 14 & 15 & 16 & 17 & 18 & 19 \\
9 & 10 & 10 & 12 & 13 & 14 & 15 & 18 & 19 & 19 \\
20 & 21 & 22 & 23 & 24 & 25 & 26 & 27 & 28 & 29 \\
20 & 21 & 22 & 23 & 24 & 25 & 27 & 28 & 29 & 29 \\
30 & 31 & 32 & 33 & 34 & 35 & 36 & 37 & 38 & 39 \\
31 & 32 & 33 & 35 & 37 & 38 & 39 & 40 & 40 & 42 \\
\end{tabular}
\end{center}

\subsection{Dyadic capacity}
For a compact nonpolar $p$-adic set $A$, let $V_p(A)$ denote the minimum logarithmic energy
\[
       V_p(A)=\inf_\mu\iint v_p(x-y)\,d\mu(x)d\mu(y),
\]
where $\mu$ ranges over nonatomic probability measures on $A$.  Equivalently,
$V_p(A)=\lim_{n\to\infty}w_A(n)/n$ for the sets considered here.

The odd branch has $V_2(H)=5/2$.  Scaling by $2$ adds one to every within-component pair valuation, so $V_2(2E)=V_2(E)+1$.  If a probability measure gives mass $t$ to $2E$ and mass $1-t$ to $H$, the cross-energy is zero and the least possible energy at that mass split is
\[
       t^2(V_2(E)+1)+(1-t)^2\frac52.
\]

\begin{theorem}[Closed-form dyadic capacity]\label{thm:dyadic-capacity}
The dyadic Robin constant of the $3$-smooth semigroup is
\begin{equation}\label{eq:V2}
       V_2(\calS)=\frac{\sqrt{11}-1}{2}=1.158312395\ldots.
\end{equation}
The equilibrium mass of the odd branch is
\begin{equation}\label{eq:qodd}
       q_2=\frac{\sqrt{11}-1}{5}=0.463324958\ldots,
\end{equation}
and the mass of the even branch is $(6-\sqrt{11})/5$.
\end{theorem}
\begin{proof}
Write $V=V_2(E)$.  Minimizing over $t$ gives the harmonic-combination equation
\[
       V=\frac{(V+1)(5/2)}{V+1+5/2}
        =\frac{5(V+1)}{2V+7}.
\]
Thus $2V^2+2V-5=0$, and positivity selects \eqref{eq:V2}.  At the minimum, the mass on $H$ is
\[
       \frac{V+1}{V+1+5/2}=\frac{\sqrt{11}-1}{5}.
\]
\end{proof}

The finite recurrence converges rapidly to this constant.
\begin{center}
\begin{tabular}{rrrr}
\toprule
$n$ & $w_2(n)$ & $w_2(n)/n$ & error from $V_2$\\
\midrule
8 & 7 & 0.875000000 & -2.833e-01 \
16 & 15 & 0.937500000 & -2.208e-01 \
32 & 33 & 1.031250000 & -1.271e-01 \
64 & 71 & 1.109375000 & -4.894e-02 \
128 & 143 & 1.117187500 & -4.112e-02 \
256 & 291 & 1.136718750 & -2.159e-02 \
512 & 586 & 1.144531250 & -1.378e-02 \
\bottomrule
\end{tabular}
\end{center}

\section{A closed formula for the generalized factorial}

Combining the three local cases gives the following explicit expression.

\begin{theorem}[Complete local factorization]\label{thm:global-factorization}
For every $n\ge0$,
\begin{equation}\label{eq:factorial-full}
 n!_{\calS}
 =2^{w_2(n)}3^{v_3(n!)}
  \prod_{p\ge5}
  p^{\displaystyle c_p\floor{n/d_p}
       +\sum_{j\ge1}\floor{n/(d_pp^j)}}.
\end{equation}
Only primes with $d_p\le n$ occur, so the product is finite.  The quantities $d_p,c_p$ are those of \eqref{eq:dp-cp}, and $w_2(n)$ is given by \eqref{eq:F-recurrence}.
\end{theorem}

For orientation, the first values are shown by prime factorization.
\begin{center}
\begin{tabular}{r l}
\toprule
$n$ & $n!_{\calS}$\\
\midrule
0 & $1$ \
1 & $1$ \
2 & $2$ \
3 & $2\,3$ \
4 & $2^{3}\,3\,5$ \
5 & $2^{4}\,3\,5$ \
6 & $2^{4}\,3^{2}\,5\,7$ \
7 & $2^{5}\,3^{2}\,5\,7$ \
8 & $2^{7}\,3^{2}\,5^{2}\,7$ \
9 & $2^{9}\,3^{4}\,5^{2}\,7$ \
10 & $2^{9}\,3^{4}\,5^{2}\,7\,11$ \
11 & $2^{10}\,3^{4}\,5^{2}\,7\,11\,23$ \
12 & $2^{10}\,3^{5}\,5^{3}\,7^{2}\,11\,13\,23$ \
13 & $2^{12}\,3^{5}\,5^{3}\,7^{2}\,11\,13\,23$ \
14 & $2^{13}\,3^{5}\,5^{3}\,7^{2}\,11\,13\,23$ \
15 & $2^{14}\,3^{6}\,5^{3}\,7^{2}\,11\,13\,23$ \
\bottomrule
\end{tabular}
\end{center}

\section{Subquadratic growth from a gcd theorem}

The odd-prime formula interacts unexpectedly well with the common-divisor sequence
\[
       G_d=\gcd(2^d-1,3^d-1).
\]
If $d_p=d$, then
\[
       c_p=v_p(G_d).
\]
The theorem of Bugeaud, Corvaja, and Zannier states that, because $2$ and $3$ are multiplicatively independent, for every $\varepsilon>0$,
\begin{equation}\label{eq:BCZ}
       \log G_d\le\varepsilon d
\end{equation}
for all sufficiently large $d$.

\begin{theorem}[Subquadratic generalized factorial]\label{thm:subquadratic}
For $\calS=\{2^a3^b\}$,
\begin{equation}\label{eq:subquad}
       \log(n!_{\calS})=o(n^2).
\end{equation}
More explicitly,
\begin{equation}\label{eq:subquad-bound}
 \log(n!_{\calS})
 \le O(n)+2n\sum_{d\le n}\frac{\log G_d}{d}.
\end{equation}
\end{theorem}
\begin{proof}
The contributions of $p=2$ and $p=3$ are $O(n)$.  For $p\ge5$, split \eqref{eq:odd-wp} into its first term and its lifting tail.  Grouping primes by $d_p=d$ gives
\[
 \sum_{p\ge5}c_p\floor{\frac n{d_p}}\log p
 \le \sum_{d\le n}\frac nd
       \sum_{p:d_p=d}c_p\log p
 \le n\sum_{d\le n}\frac{\log G_d}{d}.
\]
For the tail,
\[
 \sum_{j\ge1}\floor{\frac n{d_pp^j}}
 \le\frac{n}{d_p(p-1)}.
\]
Since $c_p\ge1$ and $1/(p-1)\le1$,
\[
 \sum_{p:d_p=d}\frac{\log p}{p-1}
 \le\sum_{p:d_p=d}c_p\log p
 \le\log G_d.
\]
This proves \eqref{eq:subquad-bound}.  Given $\varepsilon>0$, use \eqref{eq:BCZ} for all $d\ge d_0$.  The finitely many smaller $d$ contribute $O(n)$, while the remaining sum contributes at most $2\varepsilon n^2$.  Since $\varepsilon$ is arbitrary, \eqref{eq:subquad} follows.
\end{proof}

\begin{remark}[Why this estimate matters]
A generic denominator built from $n+1$ integers of exponential height can have logarithm of order $n^2$ or larger.  The theorem says that the intrinsic local denominator scale of the $3$-smooth semigroup is strictly smaller than every positive multiple of $n^2$.  Thus the local arithmetic, by itself, is small enough for a height-$e^{cn}$ interpolation contradiction.  The obstruction found below is not excessive local factorial growth; it is the impossibility of realizing all local optima with the same high nodes.
\end{remark}

\section{The first exact synchronization defect}

A sequence is a \emph{simultaneous $p$-ordering prefix} if it is a $p$-ordering prefix for every prime $p$.  For such a sequence,
\[
       \abs{\prod_{i<n}(a_n-a_i)}=n!_{\calS}.
\]
The natural small nodes initially synchronize perfectly.

\begin{proposition}[The sixth-node defect]\label{prop:sixth-node}
The chain
\[
       1,2,3,4,6
\]
is a simultaneous $p$-ordering through degree $4$:
\[
\begin{aligned}
  2-1&=1=1!_{\calS},\\
  (3-1)(3-2)&=2=2!_{\calS},\\
  (4-1)(4-2)(4-3)&=6=3!_{\calS},\\
  (6-1)(6-2)(6-3)(6-4)&=120=4!_{\calS}.
\end{aligned}
\]
It cannot be extended by a sixth node.  Indeed $5!_{\calS}=240$, whereas every new $t\in\calS$ satisfies
\[
  \abs{(t-1)(t-2)(t-3)(t-4)(t-6)}>240.
\]
The smallest continuation is $t=8$, for which the product is $1680=7\cdot240$.
\end{proposition}
\begin{proof}
The displayed products and the first generalized factorials follow from \cref{thm:global-factorization}.  A new positive $3$-smooth integer below $8$ does not exist: the elements up to $8$ are $1,2,3,4,6,8$.  For $t>6$ the polynomial
$(t-1)(t-2)(t-3)(t-4)(t-6)$ is strictly increasing, and at $t=8$ it equals $1680$.  Therefore no new semigroup node produces the required product $240$.
\end{proof}

\begin{remark}
This proposition concerns the natural chain; it does not assert that every possible six-node simultaneous ordering is impossible.  Its importance is diagnostic.  The new factor $7$ is not visible in any one local construction in advance.  It is a global synchronization cost created by forcing one integer node to satisfy all local choices simultaneously.
\end{remark}

\section{The scalar interpolation bridge}

\subsection{Arithmetic quantization of a divided difference}
Suppose $a_0,\dots,a_n$ is a simultaneous $p$-ordering prefix.  The normalized Newton polynomials
\[
 B_k(t)=\frac{\prod_{i<k}(t-a_i)}{\prod_{i<k}(a_k-a_i)}
\]
are integer-valued on $\calS$.  Their evaluation matrix at the nodes is unit lower triangular over $\Z$.  Therefore, for any function $f:\calS\to\Z$, the interpolating polynomial through the first $n+1$ values has integral Newton coefficients.  In particular,
\begin{equation}\label{eq:quantization}
       n!_{\calS}\,[a_0,\dots,a_n]f\in\Z.
\end{equation}

For $f(t)=t^x$, the generalized mean-value theorem gives, for some $\xi$ between the least and greatest node,
\begin{equation}\label{eq:meanvalue}
 [a_0,\dots,a_n]t^x
   =\frac{f^{(n)}(\xi)}{n!}
   =\binom xn\xi^{x-n}.
\end{equation}
If $x$ is not an integer and $n>x$, this number is nonzero.

\begin{theorem}[Conditional scalar bridge]\label{thm:scalar-bridge}
Assume there are integers $n_j\to\infty$ and simultaneous $p$-ordering prefixes
\[
       A_j=\{a_{j,0},\dots,a_{j,n_j}\}\subset\calS
\]
such that, with $T_j=\min A_j$,
\[
       \log T_j\ge c n_j
\]
for some fixed $c>0$.  Then the Alaoglu--Erd\H{o}s implication \eqref{eq:AE} holds.
\end{theorem}
\begin{proof}
Assume a nonintegral solution $x$.  For $j$ large, $n_j>x$.  By \eqref{eq:quantization} and \eqref{eq:meanvalue},
\[
  1\le n_j!_{\calS}\abs{\binom{x}{n_j}}T_j^{x-n_j}.
\]
For fixed nonintegral $x$,
\[
   \abs{\binom xn}
   =\frac{\Gamma(n-x)}{\abs{\Gamma(-x)}\Gamma(n+1)}
   =n^{-x-1+o(1)}.
\]
Taking logarithms and using \cref{thm:subquadratic}, the right side has logarithm
\[
      o(n_j^2)-(n_j-x)c n_j+O(\log n_j)\longrightarrow-\infty,
\]
a contradiction.
\end{proof}

The theorem is close enough to a proof that it is tempting to devote all effort to constructing high simultaneous orderings.  The next section shows that this cannot work.

\section{The two-prime synchronization obstruction}

\subsection{Incompatible equilibrium masses}
At $p=2$, an asymptotically optimal ordering has odd-node proportion $q_2$ from \eqref{eq:qodd}.  At $p=3$, an ordinary $3$-ordering is asymptotically uniform modulo $3$, so the proportion of $3$-adic units is $2/3$.

For a node $2^a3^b\in\calS$,
\[
  \text{odd}\iff a=0,
  \qquad
  \text{$3$-adic unit}\iff b=0.
\]
The two classes intersect only in the node $1$.

\begin{theorem}[No high simultaneous ordering family]\label{thm:no-high-simultaneous}
There is no sequence of simultaneous $p$-ordering prefixes $A_N\subset\calS\setminus\{1\}$ with $\#A_N=N\to\infty$.
In particular, the high-node hypothesis of \cref{thm:scalar-bridge} is impossible.
\end{theorem}
\begin{proof}
For a dyadic $p$-ordering prefix, the self-similar merge in \eqref{eq:E-decomp} and the asymptotic slopes $V_2(E)+1$ and $5/2$ force the odd proportion to converge to
$q_2=(\sqrt{11}-1)/5$.  For a $3$-ordering prefix, balance modulo $3$ forces the unit proportion to converge to $2/3$.  Since $1\notin A_N$, the odd and $3$-unit subsets are disjoint, so their proportions sum to at most $1$.  But
\[
       q_2+\frac23
       =1+\frac{3\sqrt{11}-8}{15}>1,
\]
a contradiction.
\end{proof}

This theorem explains the sixth-node defect conceptually.  It also says that no amount of brute-force search will produce the family required by the scalar bridge.

\subsection{A quantitative quadratic gap}
The incompatibility can be measured in pair energy.  For an $N$-point set $A$, define
\[
   \calE_p(A)=\frac2{N^2}\sum_{s<t\in A}v_p(s-t).
\]
Let $\theta(A)$ be its odd proportion and $\phi(A)$ its $3$-adic-unit proportion.  Put
\[
   V=\frac{\sqrt{11}-1}{2},\quad
   A_2=V+1=\frac{\sqrt{11}+1}{2},\quad
   B_2=\frac52,
\]
and
\[
   q_2=\frac{A_2}{A_2+B_2}=\frac{\sqrt{11}-1}{5}.
\]
The dyadic branch decomposition implies asymptotically
\begin{equation}\label{eq:E2mass}
 \calE_2(A)\ge A_2(1-\theta)^2+B_2\theta^2+o(1)
 =V+(A_2+B_2)(\theta-q_2)^2+o(1).
\end{equation}
At $p=3$, each of the three residue balls is a scaled copy of $\Z_3$ and has conditional energy $1+V_3(\Z_3)=3/2$.  Optimally splitting the unit mass between the two unit residue classes gives
\begin{equation}\label{eq:E3mass}
 \calE_3(A)\ge\frac32\left((1-\phi)^2+\frac{\phi^2}{2}\right)+o(1)
 =\frac12+\frac94\left(\phi-\frac23\right)^2+o(1).
\end{equation}

Set
\[
 \eta=q_2-\frac13=\frac{3\sqrt{11}-8}{15},
 \quad
 a=\frac{\sqrt{11}+6}{2}\log2,
 \quad
 b=\frac94\log3,
\]
and
\begin{equation}\label{eq:kappa}
       \kappa=\eta^2\frac{ab}{a+b}=0.0236\ldots.
\end{equation}

\begin{theorem}[Quadratic synchronization gap]\label{thm:quadratic-gap}
Let $A_N\subset\calS\setminus\{1\}$ and $\#A_N=N\to\infty$.  Then
\begin{align}\label{eq:gap}
 &\log2\left(\calE_2(A_N)-V_2(\calS)\right)
 +\log3\left(\calE_3(A_N)-V_3(\calS)\right)\\
 &\hspace{8em}\ge \kappa+o(1).
\end{align}
Equivalently, the $2$- and $3$-parts of the Vandermonde product carry at least
$\frac12\kappa N^2+o(N^2)$ more logarithmic mass than the sum of their separate equilibrium minima.
\end{theorem}
\begin{proof}
Since $1\notin A_N$, one has $\theta+\phi\le1$.  Thus, with
$u=\theta-q_2$ and $v=\phi-2/3$,
\[
       u+v\le-\eta.
\]
Multiply \eqref{eq:E2mass} by $\log2$ and \eqref{eq:E3mass} by $\log3$.  The excess is at least
$au^2+bv^2+o(1)$.  The minimum of $au^2+bv^2$ subject to $u+v\le-\eta$ is attained at equality and equals $\eta^2ab/(a+b)$.
\end{proof}

\begin{remark}
The constant in \eqref{eq:kappa} is small but strictly positive.  Its positivity, not its numerical size, is the structural fact: local minimization at $2$ and at $3$ cannot be synchronized even asymptotically once the common low node $1$ is removed.
\end{remark}

\section{Why continued-fraction clusters do not repair the scalar method}

A standard archimedean idea is to choose a good rational approximation
$u\log2\approx v\log3$ and form the near-geometric chain
\begin{equation}\label{eq:near-geometric}
       s_j=2^{ju}3^{(n-j)v},\qquad 0\le j\le n.
\end{equation}
The endpoints are a pure power of $3$ and a pure power of $2$, the gcd of all nodes is $1$, and
\[
       \frac{s_{j+1}}{s_j}=\frac{2^u}{3^v}
\]
can be arbitrarily close to $1$.  At odd primes $p\ge5$, this line can even be locally ideal when $2^u3^{-v}$ generates $H_p$.

The exceptional primes impose an exact height tax.

\begin{proposition}\label{prop:line-tax}
For the nodes in \eqref{eq:near-geometric}, let
\[
 D_j=\prod_{i\ne j}(s_j-s_i),\qquad
 L=\lcm_j\abs{D_j}.
\]
Then
\[
 v_2(L)=\frac{u n(n-1)}2,
 \qquad
 v_3(L)=\frac{v n(n-1)}2.
\]
\end{proposition}
\begin{proof}
The $2$-adic valuations of the nodes are $0,u,2u,\dots,nu$.  For $i\ne j$ these valuations are distinct, hence
$v_2(s_i-s_j)=u\min(i,j)$.  The row sum at $j$ is
\[
 u\left(\sum_{i<j}i+\sum_{i>j}j\right)
 =\frac u2j(2n-j-1),
\]
whose maximum is $un(n-1)/2$.  The $3$-adic calculation is the reverse-indexed analogue.
\end{proof}

If $u\log2\sim v\log3=L_0$, then the node height is about $nL_0$, while the logarithm of the $2$- and $3$-part of the scalar interpolation denominator is about $n(n-1)L_0$.  This cancels all but one power of the archimedean decay $T^{-n}$.  The calculation recovers the same barrier as the mass theorem in a concrete family: near-geometric clustering alone can at best exploit an exponent below $1$, whereas the conjecture must handle arbitrary positive $x$.

\section{Rank-two coefficient lattices}\label{sec:rank-two}

\subsection{Why rank two is the next natural object}
The scalar leading coefficient loses one exceptional degree of freedom at $p=2$ and another at $p=3$.  The mass obstruction is therefore codimension two.  Instead of forcing one Newton coefficient to be a nonzero rational of controlled denominator, we retain the top two coefficients of an interpolant and use their lattice covolume.

Let $A=\{a_1,\dots,a_N\}\subset\Z$ with $N=n+r$.  For integer data $y=(y_1,\dots,y_N)\in\Z^N$, let $P_y$ be the unique polynomial of degree $<N$ with $P_y(a_i)=y_i$.  Write
\[
       P_y(t)=c_0+\cdots+c_{N-1}t^{N-1}.
\]
Define the top-$r$ coefficient lattice
\begin{equation}\label{eq:coefficient-lattice}
 \calC_r(A)=
 \left\{(c_n,c_{n+1},\dots,c_{n+r-1}):y\in\Z^N\right\}
 \subset\Q^r.
\end{equation}
It is a full-rank lattice.  If $r=1$, its covolume is the reciprocal of the least common denominator of the leading-coefficient functional.  For $r=2$, it records arithmetic cancellation unavailable to any single Newton chain.

\subsection{Exact computation by Smith normal form}
Let $V_A=(a_i^{j})_{1\le i\le N,\,0\le j<N}$ be the Vandermonde evaluation matrix.  The coefficient vector is $V_A^{-1}y$.  Let $B$ be the last $r$ rows of $V_A^{-1}$.  Choose a common denominator $D$ so that $M=DB$ is integral.  The lattice generated by the columns of $M$ has covolume equal to the gcd of all nonzero $r\times r$ minors of $M$.  Hence
\begin{equation}\label{eq:covol-snf}
   \covol(\calC_r(A))
   =\frac{\gcd\{\abs{\det M_I}: I\subset\{1,\dots,N\},\ \#I=r\}}{D^r}.
\end{equation}
This is an exact integer computation.  Hermite or Smith normal form avoids rational instability and is suitable for certification in Lean.

\subsection{The analytic determinant}
Assume \eqref{eq:AE}'s hypotheses and let $P_j$ interpolate the integer values
\[
       a\longmapsto a^{x+j},\qquad j=0,\dots,r-1,
\]
on $A$.  These are integer data because $a^x\in\Z$ on $\calS$ and $a^j\in\Z$.  Let $C_x(A)$ be the $r\times r$ matrix whose $j$th column is the top-$r$ coefficient vector of $P_j$.  Every column lies in $\calC_r(A)$, so
\begin{equation}\label{eq:lattice-lower}
       \det C_x(A)\ne0
       \quad\Longrightarrow\quad
       \abs{\det C_x(A)}\ge\covol(\calC_r(A)).
\end{equation}
For positive distinct nodes, nonvanishing follows from strict total positivity of generalized Vandermonde matrices with exponent set
\[
       0,1,\dots,n-1,x,x+1,\dots,x+r-1,
\]
provided $x$ does not collide with an integer exponent in the set.  A confluent limiting argument covers repeated limiting node shapes.

If all nodes lie in $[T,CT]$, scaling $a_i=Tz_i$ shows that
\begin{equation}\label{eq:analytic-scale}
       \det C_x(A)
       =T^{r(x-n)}\,\mathfrak D_{x,n,r}(z_1,\dots,z_N).
\end{equation}
For fixed $x,r,C$, divided-difference estimates give
\begin{equation}\label{eq:shape-bound}
       \log\abs{\mathfrak D_{x,n,r}}
       =O_{x,r,C}(n\log n).
\end{equation}
The crucial point is the exponent $r(x-n)$: rank $r$ multiplies the archimedean decay without merely repeating the same scalar denominator.

\begin{criterion}[Rank-$r$ coefficient-lattice criterion]\label{crit:rank-r}
Fix $r\ge2$ and $C>1$.  Suppose there are $n\to\infty$ and sets
\[
 A_n\subset\calS\cap[T_n,CT_n],\qquad \#A_n=n+r,
 \qquad \log T_n\gg n,
\]
such that for some $\delta>0$,
\begin{equation}\label{eq:covol-target}
  -\log\covol(\calC_r(A_n))
  \le (r-\delta)n\log T_n+o(n\log T_n).
\end{equation}
Then \eqref{eq:AE} holds.
\end{criterion}
\begin{proof}
For a nonintegral solution $x$, choose $n>x$ and combine
\eqref{eq:lattice-lower}--\eqref{eq:shape-bound}.  The logarithm of the ratio of the analytic upper bound to the arithmetic lower bound is at most
\[
  r(x-n)\log T_n+(r-\delta)n\log T_n+o(n\log T_n)
  =-\delta n\log T_n+O_x(\log T_n)+o(n\log T_n),
\]
which tends to $-\infty$.  This contradicts nonvanishing.
\end{proof}

\begin{conjecture}[Rank-two synchronization conjecture]\label{conj:rank-two}
There are fixed $C>1$ and $\delta>0$ and infinitely many sets
$A_n\subset\calS\cap[T_n,CT_n]$, $\#A_n=n+2$, with $\log T_n\asymp n$, for which \eqref{eq:covol-target} holds with $r=2$.
\end{conjecture}

This is now the central computational question.  Unlike ``find a miraculous determinant,'' it has a precise input, an exact Smith-normal-form output, and a threshold dictated by a proved analytic inequality.  A counterexample to the conjecture would also be informative: it would show that the two-prime obstruction persists in exterior degree two and would tell us to raise $r$ or change the node geometry.

\subsection{Availability of high fixed-width node sets}
There is no scarcity obstruction.  For a fixed width $C\ge9$, the interval $[T,CT]$ contains $\gg\log T$ elements of $\calS$ along suitable, and in fact all sufficiently regular, heights.  Indeed, in logarithmic coordinates it is the strip
\[
  \log T\le a\log2+b\log3\le\log T+\log C,
\]
and a width at least $2\log3$ supplies at least two admissible $b$-values for a linear proportion of the possible $a$-values.  Thus $n+2$ nodes with $\log T\asymp n$ are easy to obtain.  The entire problem is their coefficient-lattice covolume.

\section{Other directions tested in this iteration}

This section records attempted routes that did not reach a proof.  Recording failures prevents repeated effort and clarifies which missing estimate is genuinely new.

\subsection{A single high divided difference}
For arbitrary integer nodes $A=\{a_0,\dots,a_n\}$ and integer data, the leading divided difference has denominator dividing
\[
       L(A)=\lcm_i\prod_{j\ne i}\abs{a_i-a_j}.
\]
The mean-value theorem makes it tempting to seek $L(A)T^{x-n}<1$.  At high $3$-smooth nodes, however, $p=2$ and $p=3$ already force almost $n$ powers of the height into $L(A)$.  The near-geometric calculation in \cref{prop:line-tax} is exact, and the mass theorem shows the phenomenon is structural.  The scalar denominator therefore has the wrong rank.

\subsection{Vanishing-moment integer measures in the integer variable}
One may choose integer coefficients $c_i$ with
\[
       \sum_i c_i a_i^j=0,\qquad 0\le j<n,
\]
so that $\sum_i c_i a_i^x$ is an integer and a Taylor remainder supplies an upper bound.  Siegel's lemma, however, makes the coefficient height depend on powers $a_i^j$.  With only $O(\log T)$ semigroup nodes in a fixed multiplicative interval, the coefficient loss overwhelms the Taylor gain.

\subsection{Vanishing moments in exponent space}
Writing $a_i=2^{u_i}3^{v_i}$ replaces the huge moment matrix by integer monomials $u_i^r v_i^s$.  If all moments of total degree $<n$ vanish, then moments of
$u\log2+v\log3$ vanish as well.  This is a genuine improvement: coefficient heights become polynomial in the logarithmic height.  The cost is that there are $\asymp n^2/2$ bivariate moment constraints.  A strip of fixed logarithmic width has only $O(\log T)$ lattice points, while a region with $O((\log T)^2)$ points is too wide for a Taylor remainder at an arbitrary exponent $x$.  Direct parameter optimization reaches only a small bounded range of $x$ and does not scale to the conjecture.  A Pad\'e version remains worth exploring, especially after reducing to the fractional part of $x$.

\subsection{Congruence transfer}
The map
\[
       2^a3^b\longmapsto m^a n^b
\]
is a monoid homomorphism and is order-preserving in the archimedean absolute value.  If it also preserved enough congruences, it would extend to maps between the local closures and force strong restrictions on $m,n$.  No such congruence preservation follows merely from real exponentiation.  In particular,
$2^a3^b\equiv2^c3^d\pmod q$ need not imply
$m^a n^b\equiv m^c n^d\pmod q$.  This route needs an additional arithmetic mechanism, not just more modular computation.

\subsection{Comparison with the output semigroup}
The output semigroup
\[
       \calS_{m,n}=\{m^a n^b:a,b\ge0\}
\]
is multiplicatively rank two and has logarithmic geometry obtained from that of $\calS$ by the scale factor $x$.  Its own generalized factorial also has subquadratic logarithm, by the same gcd theorem applied to $m,n$.  A direct monotonicity comparison between $n!_{\calS}$ and $n!_{\calS_{m,n}}$ would be powerful, but generalized factorials are local objects and real power maps do not preserve $p$-adic differences.  The comparison remains a potentially useful source of invariants, but no valid monotonicity inequality was found.

\subsection{Six exponentials amplification}
If one could prove that both $2^{x^2}$ and $3^{x^2}$ are algebraic, then the six exponentials theorem applied to the columns $1,x,x^2$ would rule out a transcendental $x$.  The identities
$2^{x^2}=m^x$ and $3^{x^2}=n^x$ show exactly what is missing: integrality of $m,n$ does not imply algebraicity of their $x$th powers.  This observation gives useful partial cases when the outputs remain inside a multiplicative group already known to be stable under $x$, but it does not settle the general problem.

\section{Computational program}

\subsection{Exact dyadic recurrence}
The following pseudocode computes $w_2(n)$ without enumerating $3$-smooth numbers.
\begin{lstlisting}[language=Python]
def vp_factorial(n, p):
    ans, q = 0, p
    while q <= n:
        ans += n // q
        q *= p
    return ans

def h(n):
    return n + n // 2 + vp_factorial(n, 2)

FH[0] = 0
for N in range(1, Nmax + 2):
    FH[N] = FH[N-1] + h(N-1)

F[0] = F[1] = 0
for N in range(2, Nmax + 2):
    F[N] = min(F[K] + K*(K-1)//2 + FH[N-K]
               for K in range(N))

w2[n] = F[n+1] - F[n]
\end{lstlisting}
Every operation is on natural numbers.  This is an unusually friendly formalization target.

\subsection{Odd-prime profile}
For each prime $p\ge5$ relevant to degree $n$:
\begin{enumerate}
\item compute $d_p=\lcm(\ord_p(2),\ord_p(3))$;
\item compute $c_p=\min(v_p(2^{d_p}-1),v_p(3^{d_p}-1))$;
\item use \eqref{eq:odd-wp}.
\end{enumerate}
A finite prime list can be generated by factoring $G_d$ for $d\le n$, because $d_p\le n$ implies $p\mid G_{d_p}$.

\subsection{Rank-two covolume search}
For a candidate node set $A$:
\begin{enumerate}
\item form the integer Vandermonde matrix $V_A$;
\item compute the last two rows of $V_A^{-1}$ exactly;
\item clear denominators by $D$;
\item compute the gcd of all $2\times2$ minors, or equivalently the Smith normal form of the resulting $2\times N$ integer matrix;
\item recover $\covol(\calC_2(A))$ from \eqref{eq:covol-snf};
\item score the set by
\[
  \sigma(A)=\frac{-\log\covol(\calC_2(A))}{n\log\min A}.
\]
The target in \cref{conj:rank-two} is a liminf strictly below $2$.
\end{enumerate}

Candidate generation should not be uniform over all subsets.  Useful families include:
\begin{itemize}
\item all $3$-smooth nodes in $[T,CT]$, followed by greedy deletion minimizing the covolume score;
\item exponent points in a logarithmic strip, selected by simulated annealing with exact Smith-normal-form rescoring;
\item unions of two near-geometric chains, one dyadically favorable and one triadically favorable;
\item sets whose exponent pairs form a digital net modulo the quotient lattices
$\{(u,v):2^u3^v\equiv1\pmod{p^k}\}$ for the small odd primes.
\end{itemize}

\section{Lean formalization blueprint}

The local results divide into small modules with minimal analytic dependencies.

\subsection{Module A: semigroup and residue images}
Define
\[
  \texttt{smooth23} = \{n:\N\mid \exists a,b,\ n=2^a3^b\}.
\]
Formalize the valuation identities
\[
 v_2(2^a3^b-2^c3^d)=\min(a,c)
 \quad\text{when }a\ne c,
\]
and the analogous identity at $3$ when $b\ne d$.
Prove surjectivity of the image modulo $3^k$ using the primitive-root property of $2$.

\subsection{Module B: procyclic local ordering}
Represent a finite quotient by a cyclic group of size $r(k)$.  Prove that the sequence of powers of a generator has occupancy counts differing by at most one and insertion cost
$\sum_k\floor{n/r(k)}$.  Instantiate it with the subgroup generated by $2$ and $3$ modulo $p^k$.

\subsection{Module C: lifting formula}
Use the standard lifting theorem for multiplicative orders at odd primes, or use $p$-adic logarithms if available, to prove
$r_p(k)=d_pp^{\max(0,k-c_p)}$.  The finite-order version can be formalized without a full $p$-adic logarithm library by applying LTE to $2^{d_p}$ and $3^{d_p}$.

\subsection{Module D: dyadic dynamic program}
Formalize the occupied-tree decomposition $E=H\sqcup2E$ at the level of finite residues.  Define $F(N)$ by well-founded recursion using \eqref{eq:F-recurrence}.  Prove by induction that it equals the minimum pair energy modulo a sufficiently deep power of $2$.  The capacity limit can initially be separated from the finite exact theorem; the recurrence and all tabulated values require only discrete arithmetic.

\subsection{Module E: sixth-node certificate}
This is the fastest end-to-end formalization target.  Prove the first five factorial values, evaluate the four products, enumerate the $3$-smooth integers below $8$, and show monotonicity of the product polynomial for $t\ge8$.  The resulting theorem is small, exact, and exercises all definitions needed later.

\subsection{Module F: coefficient lattice}
Use an integer Vandermonde matrix and define the lattice generated by the last $r$ rows of its rational inverse.  Prove \eqref{eq:covol-snf} from determinantal divisors.  This module is independent of the transcendence problem and may be reusable elsewhere in integer-valued interpolation.

\section{Prioritized next steps}

\begin{enumerate}
\item \textbf{Run the rank-two search.}  Compute $\sigma(A)$ for all manageable high strips and identify whether it trends below $2$.  This is the single most informative experiment.
\item \textbf{Prove a finite-$N$ form of the quadratic gap.}  The asymptotic theorem is enough conceptually, but an explicit inequality with an $O(N\log N)$ remainder would support certified pruning in the search.
\item \textbf{Classify small simultaneous prefixes.}  The natural chain fails at the sixth node.  An exhaustive, certificate-producing search should determine whether any six- or seven-node simultaneous prefixes exist and whether they are unique up to obvious symmetries.
\item \textbf{Study exterior-degree local trees.}  For $r=2$, derive a residue-tree recurrence for the determinantal divisor of the top-two coefficient lattice.  A self-similar dyadic formula analogous to \eqref{eq:F-recurrence} may exist.
\item \textbf{Formalize the finite core in Lean.}  Modules A, B, D, and E are independent of deep transcendence theorems.  Formalizing them now would prevent silent indexing errors and provide a verified computational kernel.
\item \textbf{Seek an averaged strengthening of BCZ tailored to $d_p$.}  The proof of \cref{thm:subquadratic} uses only the pointwise $o(d)$ estimate.  Any quantitative average bound for $\log G_d$ would sharpen the error term in coefficient-lattice thresholds.
\end{enumerate}

\section{Conclusion}

The $p$-ordering interface is now no longer a vague analogy.  It yields a complete local formula, an exact dyadic recursion, the closed constant $(\sqrt{11}-1)/2$, and a global subquadratic factorial theorem.  Those results show that a Newton interpolation proof would have enough archimedean room if the local orderings could be synchronized at high nodes.

They cannot be synchronized in rank one.  The odd mass required by the dyadic equilibrium and the unit mass required by the triadic equilibrium are incompatible away from $1$, and the incompatibility has a positive quadratic energy cost.  This converts a repeated computational failure into a theorem and prevents further investment in the wrong object.

The resulting next target is precise: determine the asymptotic covolume of the top-two coefficient lattice for high fixed-width blocks of $3$-smooth integers.  Its arithmetic is exactly computable by Smith normal form; its analytic determinant has the required scaling; and a strict saving below exponent $2$ would prove the Alaoglu--Erd\H{o}s conjecture.  Whether that saving exists is now a sharply posed finite computation and a plausible formalization project, rather than an undefined search for a miraculous identity.

\appendix
\section{Derivation of the branch-energy formulas}

For completeness, this appendix gives the elementary energy calculations behind \cref{thm:dyadic-capacity,thm:quadratic-gap}.

\subsection{A full $p$-adic ball}
For Haar measure on $\Z_p$, two independent points agree modulo $p^k$ with probability $p^{-k}$.  Therefore
\[
       V_p(\Z_p)=\sum_{k\ge1}p^{-k}=\frac1{p-1}.
\]
A ball $a+p^r\Z_p$ has energy $r+1/(p-1)$, because scaling every difference by $p^r$ adds $r$.

\subsection{The odd dyadic group}
The two components of $H$ are the balls $1+8\Z_2$ and $3+8\Z_2$.  They have energy $4$ individually.  A cross difference has valuation $1$.  Symmetry assigns mass $1/2$ to each, giving
\[
       V_2(H)=\frac14\cdot4+\frac14\cdot4
             +2\cdot\frac14\cdot1=\frac52.
\]

\subsection{Mass-constrained dyadic energy}
If odd mass is $\theta$, the even mass is $1-\theta$.  Conditional minimization gives
\[
 I_2(\theta)=(1-\theta)^2(V+1)+\theta^2\frac52.
\]
The minimizer is $q_2=(V+1)/(V+7/2)$, and completing the square gives
\[
 I_2(\theta)=V+\left(V+\frac72\right)(\theta-q_2)^2.
\]
Since $V+7/2=(\sqrt{11}+6)/2$, this is \eqref{eq:E2mass}.

\subsection{Mass-constrained triadic energy}
The three residue balls modulo $3$ each have conditional energy $3/2$.  If total unit mass is $\phi$, convexity splits it equally between the two unit balls.  Hence
\[
 I_3(\phi)=\frac32\left((1-\phi)^2+2(\phi/2)^2\right)
 =\frac12+\frac94(\phi-2/3)^2.
\]

\section{A regular-basis interpretation}

Bhargava's factorial can be described through the characteristic ideals of
\[
       \Int(\calS,\Z)=\{P\in\Q[t]:P(\calS)\subset\Z\}.
\]
For each degree $n$, the set of leading coefficients of degree-at-most-$n$ integer-valued polynomials is the fractional ideal $(1/n!_{\calS})\Z$.  A simultaneous $p$-ordering produces a Newton regular basis realizing these characteristic ideals with one sequence of nodes.  The obstruction theorem says that the local characteristic ideals exist and are small, but a high regular basis cannot be realized by one evaluation flag.  The rank-two coefficient lattice retains two characteristic directions at once and is therefore the natural exterior-power replacement.

\section{Bibliographic notes}

The original conjectural context comes from Alaoglu and Erd\H{o}s's work on highly composite and related integers.  The $p$-ordering and generalized-factorial framework is due to Bhargava.  The common-divisor estimate used in \cref{thm:subquadratic} is due to Bugeaud, Corvaja, and Zannier and is an application of the Subspace Theorem.  Standard background on four and six exponentials, linear forms in logarithms, and integer-valued polynomials can be found in the references below.

\begin{thebibliography}{99}

\bibitem{AE1944}
L. Alaoglu and P. Erd\H{o}s,
\emph{On highly composite and similar numbers},
Transactions of the American Mathematical Society \textbf{56} (1944), 448--469.

\bibitem{Bhargava1997}
M. Bhargava,
\emph{$P$-orderings and polynomial functions on arbitrary subsets of Dedekind rings},
Journal f\"ur die reine und angewandte Mathematik \textbf{490} (1997), 101--127.

\bibitem{BCZ2003}
Y. Bugeaud, P. Corvaja, and U. Zannier,
\emph{An upper bound for the G.C.D. of $a^n-1$ and $b^n-1$},
Mathematische Zeitschrift \textbf{243} (2003), 79--84.

\bibitem{CahenChabert}
P.-J. Cahen and J.-L. Chabert,
\emph{Integer-Valued Polynomials},
Mathematical Surveys and Monographs 48, American Mathematical Society, 1997.

\bibitem{LangTrans}
S. Lang,
\emph{Introduction to Transcendental Numbers},
Addison--Wesley, 1966.

\bibitem{Waldschmidt}
M. Waldschmidt,
\emph{Diophantine Approximation on Linear Algebraic Groups},
Grundlehren der mathematischen Wissenschaften 326, Springer, 2000.

\bibitem{Rumely}
R. S. Rumely,
\emph{Capacity Theory on Algebraic Curves},
Lecture Notes in Mathematics 1378, Springer, 1989.

\end{thebibliography}

\end{document}
